% ═══════════════════════════════════════════════════════════════
% 第二章：制度背景与文献综述
% ═══════════════════════════════════════════════════════════════

\section{制度背景与文献综述}

\subsection{中国上市公司监管问询制度}

监管问询函（又称"问询函"或"关注函"）是中国证券交易所对上市公司信息披露进行事中事后监管的核心工具之一。与行政处罚不同，问询函不具有法律强制力，但其公开披露的特征使其能够借助市场力量形成对上市公司的监督约束机制\citep{chen2019supervision}。

中国交易所问询函制度沿革可分为三个主要阶段。第一阶段为探索期（2014—2016年），上海证券交易所和深圳证券交易所开始对上市公司年报进行"分类监管"，对存在疑问的公司发出年报问询函。第二阶段为推广期（2017—2019年），问询函数量快速上升，问询范围从年报延伸至并购重组、关联交易、媒体报道等多个领域。第三阶段为常态化期（2020年至今），在注册制改革的推动下，问询函成为"事前审核向事中事后监管转变"的核心工具。2023年全年，两大交易所共发出各类问询函超过1500份，覆盖超过三分之一的A股上市公司。

从功能角度看，问询函具有三重信息传递效应。第一，直接效应：问询函直接指向公司的特定信息披露问题，要求公司补充说明、更正或披露更多信息。第二，信号效应：问询函的公开向市场传递了该公司可能存在信息披露缺陷的信号，已有研究发现被问询公司在问询函公告日附近会产生显著的负面超额收益\citep{li2019inquiry}。第三，威慑效应：问询函的存在对潜在违规公司形成威慑，有助于提升市场整体信息披露质量\citep{chen2019effective}。

\subsection{信息披露、软信息与公司风险评估}

公司风险评估是金融经济学和会计学的核心议题。传统上，研究者主要依赖结构化财务指标来预测公司风险和困境。Altman（1968）提出的Z-Score模型利用多变量判别分析预测破产风险，是这一领域的开创性工作\citep{altman1968financial}。Beneish（1999）提出的M-Score模型则聚焦盈余操纵识别，利用八个财务指标综合判断公司粉饰报表的可能性\citep{beneish1999detection}。这些方法虽然可解释性强、计算简便，但其局限也日益显现：仅依赖结构化财务数据，难以捕捉公告文本、管理层讨论等非结构化信息中所蕴含的风险信号。

近年来，文本分析在金融经济学中的应用快速发展。在方法论层面，已有研究利用词典法（如Loughran-McDonald情感词典）、机器学习（如支持向量机、随机森林）和预训练语言模型（如BERT及其金融领域变体FinBERT）提取文本特征，在盈余预测、股价波动预测和信用风险评估等任务中取得了超越纯结构化模型的性能\citep{loughran2011liability,huang2023finbert}。Liberti和Petersen（2019）在《经济文献杂志》（JEL）发表的综述中将信息划分为"硬信息"（hard information，如财务指标）和"软信息"（soft information，如文本描述、管理层语调），系统论证了软信息在信贷决策、投资分析和风险定价中的独特价值\citep{liberti2019information}。

在信息披露质量研究领域，Healy和Palepu（2001）的综述奠定了理论基础，指出高质量的信息披露通过降低信息不对称减少代理成本、提升资本市场效率\citep{healy2001information}。近期研究进一步发现，管理层讨论与分析（MD\&A）的语调\citep{davis2015managerial}、可读性\citep{li2008readability}和前瞻性陈述密度\citep{muslu2015forward}均包含预测公司未来业绩的增量信息。

\subsection{监管问询的经济后果研究}

监管问询的经济后果已成为中国金融学研究的热点领域。陈运森等（2019）系统检验了交易所一线监管的有效性，发现年报问询函能显著改善公司的信息披露质量，并产生积极的治理效应\citep{chen2019supervision}。李晓溪等（2019）发现年报问询函能够提升管理层业绩预告的准确性，减少管理层与市场之间的信息不对称\citep{li2019inquiry}。李晓溪等（2020）进一步发现问询函在并购重组中发挥了信息甄别功能，有助于抑制并购溢价和改善并购绩效\citep{li2020ma}。

在投资者反应方面，已有研究一致发现问询函公告产生显著的负向市场反应\citep{chen2019effective}，且问询函的问题数量和涉及金额越大，市场反应越强烈。在分析师反应方面，问询函公告后分析师对目标公司的盈利预测准确性有所提高\citep{wang2019analyst}。在公司行为反应方面，被问询公司在收到问询函后会显著改善其信息披露质量，表现为可读性提升、语调改善和补充披露增加\citep{li2019inquiry}。

然而，现有文献主要集中在监管问询的\textit{事后效应}——即问询函发出后产生了怎样的经济和市场影响。对于\textit{事前预测}——即哪些公司特征可以提前预测监管问询——的研究尚处于起步阶段。少数学者利用Logistic回归和传统财务指标对问询概率进行了初步探索\citep{cassell2013review}，但受限于数据维度单一和方法传统，现有研究的预测精度有限，且未能有效利用非结构化的文本信息。

\subsection{研究空白与本文定位}

综上，已有研究存在以下三点不足：第一，在数据维度上，现有预测研究主要依赖财务指标，对公告文本语义、市场微观结构、公司治理特征等异构信息的融合利用不足；第二，在方法工具上，大语言模型为代表的先进文本分析技术在监管问询预测中的应用尚属空白，LLM提取的"软信息"相对于传统"硬信息"的增量预测价值有待检验；第三，在可解释性上，现有预测模型多为"黑箱"式输出，缺乏从概率分数到风险特征再到原文证据的完整可追溯归因机制，难以满足投研和风控人员的复核需求。

本文在上述三个维度上做出贡献。在数据维度，本文构建了涵盖六大类约235个特征的多元信息融合框架，将财务指标、LLM语义特征、市场特征、历史监管记录、知识图谱关联特征和时序衍生特征纳入统一分析框架。在方法维度，本文首次将大语言模型提取的文本语义风险特征用于监管问询概率预测，并系统检验其相对于传统财务指标的增量预测价值。在可解释性维度，本文构建了"SHAP归因—LLM生成报告—原文证据追溯"三级可解释机制，实现了从概率输出到具体证据的完整归因链。